\begin{proposition}[Absence de famine sur un ensemble fini]
Soit $S$ un ensemble fini de candidats continuellement admissibles. Supposons que chaque candidat ait un coût inférieur ou égal à la capacité allouable lors d'une ronde de service, qu'aucune protection (pin, lease, transaction) ne l'exclue, que le contrôleur soit exécuté indéfiniment, qu'une part positive des rondes serve la classe concernée, et que l'ordonnancement sélectionne en priorité les candidats de cette classe par âge décroissant avec remise à zéro de l'âge servi. Alors aucun candidat de $S$ ne peut être différé indéfiniment; à partir du moment où la classe est continuellement éligible, chaque candidat finit par être activé.
\end{proposition}
\begin{proof}
À chaque ronde de service, un candidat sélectionné voit son âge remis à zéro; chaque candidat continuellement admissible et non sélectionné vieillit. L'ensemble $S$ étant fini, il existe à chaque ronde un candidat d'âge maximal. La règle de priorité empêche un candidat récemment servi de passer indéfiniment devant un candidat plus ancien. Si un candidat restait indéfiniment sans service, son âge finirait par dépasser celui de tout candidat sélectionné dans la classe; l'ordonnanceur devrait alors le sélectionner, contradiction. La conclusion dépend des hypothèses de coût, d'admissibilité continue et de rondes de service; elle ne couvre pas un espace infini ni un flot d'arrivées non borné.
\end{proof}
